\scanpage{96}
plus une question pour des courbes relatives \uline{globales}, mais pour des ``éléments de courbe'' relatifs --- savoir, pour le complété formel d'une courbe relative $X/B$ le long d'une section $g$, avec le diviseur de cette section distingué dans la structure.

Je vais plutôt revenir à la situation où on part de une section de première espèce d'un $\mathcal E_{U_\eta}$ sur $\mathcal E_\eta$, \unsure{associée|associé} donc à un pt \unsure{rationnel|rationnelle} $f_\eta$ de $U_\eta$ sur $K$. On peut donc dire (indépendamment de la donnée du modèle $S$ de $K$ sur lequel $U_\eta$ pouvait déjà se définir en \deleted{\unsure{tant}} que schéma relatif élémentaire) que la dite section \deleted{de} $\mathcal E_{U_\eta}$ sur $\mathcal E_\eta$ \inserted{$(\unsure{29|23})\quad\unsure{\mathcal E f_\eta|f_\eta}=\psi:\mathcal E_\eta\longrightarrow\mathcal E_{U_\eta}$} définit, pour \unsure{tte|tt} section de $\mathcal E_K$ sur $\mathcal E_{\mathbf Q}=\Gamma_{\mathbf Q}$ ou sur un sous-groupe ouvert \inserted{$\Gamma'$} de celui-ci, une section
\margin{95}
